% Pasteable draft for two tutorial sections.
% Light markup only: sectioning, enumerate, math. No Lean syntax.
% Cite keys: deMoura2021, mathlib2020, Avigad2014 (Avigad--Harrison).
% Intended audience: materials scientists who have never heard of Lean.

\section{What Lean is and why it is here}
\label{sec:what-lean-is}

A proof assistant is a computer program that checks a mathematical
argument the way a compiler checks types: every inference must follow
from stated definitions and previously accepted lemmas, or the file is
rejected.
Lean~4 is such a system~\cite{deMoura2021}.
Its companion library Mathlib is a large, continuously checked corpus
of algebra, analysis, and special functions~\cite{mathlib2020}.
Together they let us write the Fourier-optics identities of Yu, Lau and
Altman~\cite{Yu2019} so that the hypotheses cannot silently drift.
That is a different activity from writing a Python script.
A script evaluates numbers: it can plot a through-focal series, invert
a matrix, or report a residual.
It does not, by itself, certify that the kernel on the screen is the
kernel in the paper, that the regularized inverse is the unique
minimizer of the energy we wrote down, or that an error bound is sharp
rather than a lucky numerical observation.
Surveys of this distinction, and of why machine-checked proofs matter
in scientific computing, are given by Avigad and
Harrison~\cite{Avigad2014}.

The LEEM inverse is exactly the setting in which that distinction
bites.
The published forward model is bilinear Fourier optics (FO): intensity
is assembled from a kernel $R_{\mathrm{FO}}(\mathbf{q},\mathbf{q}',\Delta z)$
that does not factor under partial coherence.
A Python reconstruction that divides by a multiplicative
contrast-transfer function (CTF) is therefore solving a different
problem as soon as the object samples off-axis pairs
$(\mathbf{q},\mathbf{q}')$.
Even if one stays on the CTF slice $\mathbf{q}'=\mathbf{0}$, five
conventions are easy to get wrong and hard to notice in a plot.
The optics kernel uses the column wavelength $\lambda$ after
acceleration; the kinematic object phase $\phi=4\pi A/\lambda_0$ uses
the wavelength $\lambda_0$ at the sample.
The wave aberration $\chi$ is in waves, so the transfer factor is
$\mathrm{e}^{2\pi\mathrm{i}\chi}$, not $\mathrm{e}^{\mathrm{i}\chi}$ with
$\chi$ already in radians.
The CTF picture and the FO picture agree on the axes and disagree off
them; using one as if it were the other is the error Yu and
coworkers documented for rippled graphene.
Defocus has a sign: this work takes
$+\tfrac12\Delta z\,\lambda q^2$ in $\chi_S$, so that positive
$\Delta z$ adds optical path for off-axis rays; mixing a minus-sign
convention reverses overfocus and underfocus without changing a single
plot label.

What, then, does it mean to say that a ``numerical proof is good''?
It does not mean that a reconstruction looks plausible, or that a
residual is small on one dataset.
It means three things that a script can illustrate but cannot settle.
First, the regularized estimator is unique: for every frequency bin and
every Tikhonov parameter $\alpha>0$ there is exactly one minimizer,
given by a closed formula.
Second, the error of that estimator is an identity, not a fitted slope:
the difference between the estimate and the true bin amplitude is
exactly a projected discrepancy minus a bias term, over a positive
denominator.
Third, the method has physical blind spots that no amount of
regularization removes---modes outside the contrast aperture are
identically invisible; a global phase of the exit wave never appears
in the intensity; a single defocus cannot identify a general complex
object; and the weak-phase CTF has zeros inside a large enough
aperture.
Those are the claims Lean is here to lock down.
The numerical method remains the method; Lean is the certificate that
we are running the method we think we are running.

\section{What Lean certifies so the numerical method is trustworthy}
\label{sec:what-lean-certifies}

The discrete linear stage is trustworthy because the following
statements have been checked, not because a reconstruction on one stack
looked clean.
They are listed in the order a microscopist meets them.

\begin{enumerate}
\item
  \textbf{Same forward model $R_{\mathrm{FO}}$ and CTF slice.}
  The reconstruction uses the bilinear kernel of Yu~\emph{et al.}
  and the same CTF slice: $R_{\mathrm{FO}}$ and $R_{\mathrm{CTF}}$
  coincide when one frequency argument is zero.
  Stage~1 inverts that slice, not a substitute propagator.

\item
  \textbf{Unique regularized minimizer per frequency bin.}
  After a discrete Fourier transform, each spatial frequency decouples.
  For $\alpha>0$ the Tikhonov energy has a unique global minimizer,
  given in closed form by a weighted inner product of the measurements
  against the slice multipliers, divided by $\alpha$ plus the sum of
  their squared moduli.
  No other regularized solution exists for that energy.

\item
  \textbf{Exact error formula.}
  If the data equal the linear model plus a discrepancy $n$
  (instrument noise, discretization, and the quadratic FO remainder
  lumped together), the estimation error is exactly
  $\bigl(\sum_k \overline{h_k}\,n_k-\alpha x^\star\bigr)/D$ with
  $D=\alpha+\sum_k\lvert h_k\rvert^2$.
  This is an algebraic identity.
  It is not a statistical noise theorem: $n$ need not be random,
  Gaussian, or zero-mean.

\item
  \textbf{Triangle bound, sharp.}
  The modulus of that error is at most
  $\bigl(\sum_k \lvert h_k\rvert\,\lvert n_k\rvert
  +\alpha\lvert x^\star\rvert\bigr)/D$.
  The constant one in front of the bound cannot be improved uniformly;
  equality is attained on an explicit one-measurement example.

\item
  \textbf{Physical limits (aperture, gauge, $K=1$, CTF zeros).}
  Frequencies outside the contrast aperture have vanishing slice
  multipliers at every defocus, so those bins reconstruct as zero.
  Intensity is unchanged if the exit wave is multiplied by a global
  phase; the linear inverse around vacuum therefore chooses a gauge
  rather than recovering an absolute phase.
  One image never identifies a general complex pair
  $\bigl(X(\mathbf{q}),X(-\mathbf{q})\bigr)$.
  The weak-phase factor $\sin(2\pi\chi_S)$ vanishes on rings inside a
  large enough aperture of a defocused, non-aberration-corrected
  column; those holes are why one stacks at least two independent
  defoci.

\item
  \textbf{The linear method is first-order; the quadratic remainder
  licenses stage~2.}
  The multi-defocus slice map is exactly the derivative of bilinear FO
  at vacuum.
  The neglected term is the FO image of the increment itself,
  homogeneous of degree two.
  When that remainder sits above the noise, Gauss--Newton on the full
  bilinear residual is the consistent next step; when it does not,
  stage~2 cannot help.

\item
  \textbf{Cost model, not a benchmark.}
  Counting a modelled discrete Fourier transform as
  $N\log_2 N$ operations and a per-bin Gram solve as a few tens of
  arithmetic operations, the linear stage is $O(KN\log N)$ and,
  already at a $128$-pixel grid, strictly cheaper than a dense
  bilinear apply.
  Existence of a fast Fourier transform, wall-clock times, and cache
  behaviour are not proved.

\item
  \textbf{Explicitly not proved.}
  Alignment of the experimental stack, the map from lens current to
  $\Delta z$, equality of the array transform with the continuum
  Fourier transform of the paper, Gaussian or Poisson statistics and
  Cram\'er--Rao bounds, local quadratic convergence of nonlinear
  Gauss--Newton, and any claim to be fastest among all conceivable
  algorithms, remain informal.
\end{enumerate}

\paragraph{Stage~1 recipe.}
\begin{enumerate}
\item
  Align the through-focal stack by rigid shifts; this is
  preprocessing, not part of the certificate.
\item
  Subtract the spatial mean of each plane and take a discrete Fourier
  transform, so that the zero-frequency bin, which carries total flux,
  is handled separately.
\item
  For every frequency inside the aperture and every recorded defocus,
  evaluate the closed-form slice
  $h_k(\mathbf{q})=R_{\mathrm{FO}}(\mathbf{q},0,\Delta z_k)$ from the
  Yu~2019 envelopes, using column $\lambda$, $\chi$ in waves, and the
  $+\tfrac12\Delta z\,\lambda q^2$ sign.
\item
  Solve the scalar, or conjugate-pair $2\times 2$, Tikhonov problem at
  that bin with a fixed $\alpha>0$, and set the reconstruction to zero
  outside the aperture.
\item
  Inverse-transform the assembled spectrum and read amplitude and phase
  in the vacuum gauge, with the zero-frequency component real and
  positive.
\item
  Evaluate the bilinear FO residual of this reconstruction against the
  stack: if it is consistent with the noise floor, stop, because the
  quadratic remainder is undetectable; if not, the same residual
  licenses a few Gauss--Newton steps on full FO (stage~2), or, for the
  one-dimensional sinusoid of Yu~\emph{et al.}, a Jacobi--Anger fit of
  the Bessel ladder.
\end{enumerate}

The discrete linear stage is recommended because, under the Yu~2019
forward model, it is well-posed: unique for every $\alpha>0$, equipped
with an exact error identity and a sharp bound, honest about aperture,
gauge, single-image, and CTF-zero blindness, and first-order consistent
with bilinear FO.
Lean is the certificate of those statements.
The numerical method is what one runs; the certificate is why one may
trust that the method is the one the optics requires.

% Suggested bibliography entries (author--year keys as requested):
%
% [deMoura2021] L. de Moura, S. Ullrich, The Lean 4 Theorem Prover and
%   Programming Language, CADE 28, LNCS 12699 (2021) 625--635.
% [mathlib2020] The mathlib Community, The Lean mathematical library,
%   CPP 2020, ACM (2020) 367--381.
% [Avigad2014] J. Avigad, J. Harrison, Formally verified mathematics,
%   Commun. ACM 57 (4) (2014) 66--75.
% [Yu2019] K. M. Yu, K. L. W. Lau, M. S. Altman, Ultramicroscopy 200
%   (2019) 160--168.
